% getting_started.tex -- CLI-first guide from installation to first sync

\chapter{Getting Started}

This chapter guides a new user through the Multi-repo Sync CLI. The documented
entry point is \cgscli{}.

The primary lifecycle is:

\begin{enumerate}
  \item \texttt{initialise(.cgs)} $\rightarrow$ clone all repositories
        $\rightarrow$ \texttt{READY}
  \item \texttt{initialise(.gts)} $\rightarrow$ load or restore a saved
        snapshot $\rightarrow$ \texttt{READY}
  \item \texttt{pull(.cgs|.gts)} $\rightarrow$ resynchronise an existing tree
  \item \texttt{checkout} / \texttt{add} / \texttt{commit} / \texttt{push} /
        \texttt{tag} $\rightarrow$ tree-wide Git operations
  \item \texttt{freeze} $\rightarrow$ write the next \texttt{.gts} snapshot
        and update the project-local \texttt{.lgr}
\end{enumerate}

% -----------------------------------------------------------------------
\section{Prerequisites}

\begin{itemize}
  \item Python~3.11 or later.
  \item Git~2.30 or later available on \texttt{PATH}.
  \item Working Git authentication for every remote used by the workspace.
\end{itemize}

% -----------------------------------------------------------------------
\section{Installation}

Install \cgspkg{} from source with the Pixi-managed environment:

\begin{lstlisting}[language=bash]
$ git clone https://github.com/flipoyo/ComplexGitSync.git
$ cd ComplexGitSync
$ pixi install
$ pixi run cgitsync --help
\end{lstlisting}

% -----------------------------------------------------------------------
\section{Create or Check a Project Spec}

A \texttt{.cgs} file is the local authoring spec for the repository tree. It
normally contains only a project name and repository identifiers:

\begin{lstlisting}
project = "CGSil1"
repos = [
  "gitlab:CGS_test/CGSil1",
  "gitlab:CGS_test/CGSil2",
  "github:flipoyo/CGSih1",
]
\end{lstlisting}

Each identifier uses \texttt{provider:owner/repository}.  ComplexGitSync runs
\texttt{PARSE -> NORMALIZE -> VALIDATE} and supplies deterministic defaults:
\texttt{main} branches, \texttt{ssh}, \texttt{nested\_config = auto}, and the
repository name as its path.  A unique repository matching the project name
is mounted at \texttt{.}.  Advanced tables are reserved for overrides.

Canonical providers are \texttt{github}, \texttt{gitlab},
\texttt{codeberg}, and \texttt{custom}. Their known hosts are
\texttt{github.com}, \texttt{gitlab.com}, and \texttt{codeberg.org}; SSH and
HTTPS remotes are generated automatically. Custom hosting requires an explicit
\texttt{gitprovider\_url} and never guesses a host. For example,
\texttt{codeberg:GX4G/GX4G} identifies owner \texttt{GX4G} and repository
\texttt{GX4G} on Codeberg.

Copy \texttt{examples/template.cgs} when starting a specification. It combines
plain repository identifiers with one inline override. Developers can compare
it with \texttt{examples/normalized\_template.cgs}, which shows the complete
canonical document after normalization and is not intended as the usual
hand-authored form.

The common inline options are \texttt{default\_branch},
\texttt{fallback\_branch}, \texttt{access\_protocol},
\texttt{relative\_path}, and \texttt{nested\_config}. A relative path is
resolved from the parent repository. Nested config can be \texttt{auto} to
load the sole root-level \texttt{*.cgs}, \texttt{disabled} to stop discovery,
or a relative \texttt{.cgs} path to select an exact file inside the repository.

\begin{lstlisting}[language=bash]
$ pixi run cgitsync validate examples/complexgitsync.cgs
$ pixi run cgitsync view-tree examples/complexgitsync.cgs
\end{lstlisting}

A valid declared tree prints a lifecycle line such as:

\begin{lstlisting}[language=bash]
DECLARED ready=false complete=true
\end{lstlisting}

% -----------------------------------------------------------------------
\section{Initialise the Workspace}

Use \texttt{initialise} as the canonical first command. From a \texttt{.cgs}
file it clones the full repository tree, validates it, writes the runtime
snapshot under \texttt{<workspace>/.cgitsync/state/}, and leaves the tree
\texttt{READY}.

\begin{lstlisting}[language=bash]
$ pixi run cgitsync initialise examples/complexgitsync.cgs
\end{lstlisting}

The same canonical project definition can be authored directly. The
\texttt{--repo} option is repeatable:

\begin{lstlisting}[language=bash]
$ pixi run cgitsync initialise --project CGSil1 \
    --repo gitlab:CGS_test/CGSil1 \
    --repo codeberg:GX4G/GX4G
\end{lstlisting}

Use either a source file or \texttt{--project} plus \texttt{--repo}, never
both. To serialize the CLI definition without initializing a workspace, use
\texttt{create-cgs --project ... --repo ... --output FILE}.

When \texttt{--output-path} is omitted, \cgscli{} uses the workspace-level
default \texttt{CGSPATH=../..}, then derives
\texttt{CGSHOME=CGSPATH/<project-name>}. The explicit equivalent is:

\begin{lstlisting}[language=bash]
$ pixi run cgitsync initialise examples/complexgitsync.cgs --output-path "$CGSPATH"
\end{lstlisting}

From a previously generated snapshot (resolved automatically via the
project's \texttt{.lgr} register; pass the path explicitly only to override
that discovery):

\begin{lstlisting}[language=bash]
$ pixi run cgitsync initialise
\end{lstlisting}

Successful \texttt{.cgs} initialisation prints a readable operation sequence
\texttt{GT-LOAD->GT-DISCOVER->GT-VALIDATE->GT-CLONE->GT-GITIGNORE}, the
explicit workflow \texttt{load->expand->validate->clone->gitignore}, the
per-repository \texttt{git clone} action, the resolved root path, a report of
any \texttt{.gitignore} changes, and a compact tree outline. Structured log
records also put the operation code first, for example \texttt{GT-DISCOVER},
\texttt{GT-VALIDATE}, \texttt{FS-PURGE}, \texttt{GT-CLONE}, or
\texttt{GT-GITIGNORE}.
If a stale partial checkout blocks the clone step,
\texttt{initialise} fails explicitly and prints
\texttt{Try clean-init method}.  Use \texttt{clean-init} to insert a purge
phase between validation and cloning:

\begin{lstlisting}[language=bash]
$ pixi run cgitsync clean-init examples/complexgitsync.cgs
\end{lstlisting}

The explicit operation sequence is
\texttt{GT-LOAD->GT-DISCOVER->GT-VALIDATE->FS-PURGE->GT-CLONE->GT-GITIGNORE};
the matching workflow is \texttt{load->expand->validate->purge->clone->gitignore}.
The purge phase can also be run on its own:

\begin{lstlisting}[language=bash]
$ pixi run cgitsync purge examples/complexgitsync.cgs
\end{lstlisting}

\texttt{purge} removes generated clone state from \texttt{CGSHOME}: immediate
child repository directories declared directly under the project root, and
project \texttt{*.lgr} files. A persisted Git identity override at
\texttt{.cgitsync/master.toml} (see the \texttt{initialise} entry in the
Command Reference chapter for \texttt{--git-user-name}/\texttt{--git-user-email})
is workspace configuration, not clone state, and is left untouched.

% -----------------------------------------------------------------------
\section{Inspect the Runtime State}

\begin{lstlisting}[language=bash]
$ pixi run cgitsync view-tree
$ pixi run cgitsync view-tree --depth 2 --collapse parent_2
$ pixi run cgitsync status
\end{lstlisting}

If a command accepts an optional snapshot and none is provided, \cgscli{}
discovers the latest \texttt{.gts} automatically via the project's
\texttt{.lgr} register under \texttt{CGSHOME/.cgitsync/}. Each generated
\texttt{.gts} snapshot actually lives under its own content-addressed
\texttt{CGSHOME/.cgitsync/state(<hash>)\_<n>/} directory; pass an explicit
path (e.g. \texttt{--gts FILE}, or a positional source, depending on the
command) only to override that discovery.
\texttt{CGSHOME} can be set explicitly; for \texttt{initialise(.cgs)} the
default is \texttt{CGSPATH=../..}, and later commands can also discover the
workspace by walking upward from the current directory.

% -----------------------------------------------------------------------
\section{Run the Git Cycle}

Once the tree is \texttt{READY}, the Multi-repo Sync commands operate on every
repository in deterministic order:

\begin{lstlisting}[language=bash]
$ pixi run cgitsync pull
$ pixi run cgitsync checkout feature/my-branch
$ pixi run cgitsync add
$ pixi run cgitsync commit "feat: update project CGS#1"
$ pixi run cgitsync push
$ pixi run cgitsync freeze-release release-2026.05 "release 2026.05"
$ pixi run cgitsync freeze release-2026.05
$ pixi run cgitsync launch-release release-2026.05
\end{lstlisting}

\texttt{pull} walks the tree parent-first, including the project root:
\texttt{root -> parent -> leaf}.  Mutation commands such as \texttt{add},
\texttt{commit}, \texttt{push}, and \texttt{freeze} walk the opposite direction:
\texttt{leaf -> parent -> root}.
When local files block the safe pull, the CLI suggests
\texttt{cgitsync pull-force}; that recovery command is destructive and discards
local uncommitted/untracked work before force-updating the same tree.

Pass \texttt{--gts <snapshot.gts>} when you want to pin a command to an
explicit snapshot, and pass \texttt{--dry-run} to \texttt{add},
\texttt{commit}, \texttt{push}, or \texttt{freeze} to preview
the plan without mutating repositories.

% -----------------------------------------------------------------------
\section{Log Files}

Every CLI command writes a timestamped log file. On Linux the default location
is \texttt{~/.local/state/ComplexGitSync/logs/}; if
\texttt{XDG\_STATE\_HOME} is set, logs are written under that state directory
instead. The CLI prints the resolved path as \texttt{log\_file=<path>}.

% -----------------------------------------------------------------------
\section{Next Steps}

See Chapter~\ref{chap:user-guide} for the command reference, document formats,
preflight checks, and troubleshooting notes.
